Thus, the probability of observing \(Y=1\) is
\[
\Pr(Y=1 \mid X) \;=\; \Phi\bigl(X\beta\bigr),
\]
where \(\Phi(\cdot)\) is the standard normal cumulative distribution function.

The sample log-likelihood is given by
\[
\ell_Y(\beta)
\;=\;
\sum_{i=1}^n \left[
Y_{i}\,\ln \Phi\!\bigl(X_i\beta\bigr)
\;+\;
(1-Y_{i})\,\ln \Bigl(1-\Phi\!\bigl(X_i\beta\bigr)\Bigr)
\right].
\]
Maximizing this log-likelihood with respect to \(\beta\) yields the maximum likelihood estimates \(\hat{\beta}\).

\bigskip
\noindent
\textbf{Grade Inflation.} I define school-level grade inflation as the change in the predicted probability of achieving an \(A\) (or \(A^*\)) between the baseline period and the treatment year (2020), holding other covariates at a reference level \(\overline{X}\). For school \(j\), the inflation factor \(\theta_j\) is computed as
\[
\theta_j \;=\;
\Pr\bigl(Y=1 \mid \text{Year2020}=1,\text{School}=j,\overline{X}\bigr)
\;-\;
\Pr\bigl(Y=1 \mid \text{Year2020}=0,\text{School}=j,\overline{X}\bigr).
\]
In the context of my logit model, these probabilities are evaluated as
\[
\Pr(Y=1 \mid X) \;=\; \Phi\Bigl( X\hat{\beta} \Bigr).
\]

As I estimate the school level treatment effect $\{\delta_j\}_{j=1}^{4000}$ using a ridge regression. Estimates of the school level treatment effects obtained through standard maximum likelihood algorithms exhibit high variance (Appendix []). 
